% -------------------------------------------------------
%  Abstract
% -------------------------------------------------------


\شروع{وسط‌چین}
\مهم{چکیده}
\پایان{وسط‌چین}
\بدون‌تورفتگی

مسئلهٔ زمان‌بندی کارگاهی\پاورقی{Job Shop Scheduling Problem (JSSP)}
یکی از مسائل بنیادیِ بهینه‌سازی ترکیبیاتی است که در برنامه‌ریزی تولید و سامانه‌های صنعتی
کاربرد گسترده‌ای دارد.
این مسئله در ردهٔ مسائل ان‌پی-سخت قرار می‌گیرد؛ ازاین‌رو دستیابی به پاسخ بهینه،
به‌ویژه برای نمونه‌های بزرگ‌مقیاس و صنعتی، مستلزم هزینه محاسباتی قابل‌توجهی است.
در این پایان‌نامه، نسخهٔ چندهدفه‌ای از این مسئله در چارچوب یک زیرساخت رایانشی ناهمگنِ سه‌لایه
(لبه-مه-ابر) صورت‌بندی می‌شود؛ به‌گونه‌ای که علاوه بر تعیین توالی اجرای عملیات،
برای هر عملیات یک گرهٔ رایانشی نیز تخصیص می‌یابد.
اهداف مسئله عبارت‌اند از:
کمینه‌سازی زمان تکمیل کل،
کمینه‌سازی مصرف کل انرژی،
و کمینه‌سازی عدم‌قابلیت‌اطمینان سامانه (یا معادل آن، بیشینه‌سازی قابلیت اطمینان).

برای حل مسئله، چارچوب الگوریتم ژنتیک چندهدفهٔ \lr{NSGA-II}
با سه سازوکار مکمل تقویت شده است:
(۱)~مقداردهی اولیهٔ هوشمند با تلفیق بذرهای ابتکاری و تنوع تصادفی،
(۲)~انتخاب تطبیقی عملگر بر پایهٔ سازوکار راهزن چندبازویی و قاعدهٔ \lr{UCB1}،
و (۳)~به‌کارگیری جست‌وجوی محلی در قالب مؤلفهٔ میمتیک.
ارزیابی راه‌حل‌ها با شبیه‌سازی زیرساخت رایانشی لبه-مه-ابر
در شبیه‌ساز \lr{YAFS} انجام شده است و هر سه تابع هدف به‌صورت قطعی محاسبه می‌شوند.

ارزیابی تجربی روی ۱۵ نمونهٔ معیار از خانواده‌های
\lr{FT}، \lr{LA}، \lr{ORB}، \lr{ABZ}، \lr{TA}، \lr{SWV} و \lr{YN}
و در قالب ۱۰ اجرای مستقل انجام شده است.
نتایج با شاخص‌های ابرحجم و فاصله نسلی معکوس سنجیده شده
و معنی‌داری آماری آن‌ها با آزمون‌های ناپارامتری فریدمن و ویلکاکسون
(به‌همراه تصحیح هولم) بررسی شده است.
یافته‌ها نشان می‌دهد الگوریتم ژنتیک بهبودیافته در اغلب نمونه‌ها،
در مقایسه با نسخهٔ پایه و قواعد ابتکاری، جبهه‌های پارتوی باکیفیت‌تری تولید می‌کند.
همچنین مطالعهٔ حذفی بیانگر آن است که انتخاب تطبیقی عملگر و جست‌وجوی محلی،
هر دو سهم معناداری در بهبود عملکرد دارند.

\پرش‌بلند
\بدون‌تورفتگی \مهم{کلیدواژه‌ها}:
زمان‌بندی کارگاهی، بهینه‌سازی چندهدفه، الگوریتم ژنتیک، \lr{NSGA-II}، رایانش لبه‌ای-مه-ابر، ابرحجم
\صفحه‌جدید
